% arara: clean: {
% arara: --> extensions:
% arara: --> ['aux','log','out','blg','idx','ind',
% arara: --> 'toc','bbl', 'bcf','run.xml','ilg','pdf']
% arara: --> }
% arara: lualatex: {
% arara: --> shell: 1,
% arara: --> interaction: batchmode
% arara: --> }
% arara: biber
% arara: makeindex
% arara: lualatex: {
% arara: --> shell: 1,
% arara: --> interaction: batchmode
% arara: --> }
% arara: lualatex: {
% arara: --> shell: 1,
% arara: --> interaction: batchmode
% arara: --> }
% arara: clean: {
% arara: --> extensions:
% arara: --> ['aux','log','out','blg','idx','ind',
% arara: --> 'toc','bbl', 'bcf','run.xml','ilg']
% arara: --> }

\documentclass[
  fontsize=12pt,
  answers,
  addpoints,
  pagesize,
  paper=landscape,
  a5,
  DIV=calc,
  headsepline,
  footsepline,
  abstract=true,
  bibliography=totoc
]{kommaexam}

\usepackage[margin={.38cm}]{geometry}

\usepackage[brazilian]{babel}
\usepackage{csquotes}

\usepackage[useregional]{datetime2}
\usepackage[svgnames]{xcolor}
\usepackage{graphicx,svg}
\usepackage{caption,subcaption}
\usepackage[shortlabels]{enumitem}

%\usepackage{mathpazo}
\usepackage{mathtools,amssymb,amsthm}

\usepackage{fontspec}%unicode-math
\usepackage{newunicodechar}
\setmainfont{TeX Gyre Pagella}
\setmonofont{Fira Code}[Contextuals=Alternate]


\usepackage[chapter]{minted}
\newmintedfile[juliacode]{julia}{
	breaklines,
	breakanywhere,
	tabsize=4,
	samepage=false,
	showspaces=false,
	showtabs =false
}
%\usepackage{pdfpages}
\graphicspath{ {img/} }

\usepackage{background}
\usepackage{tikz}
\usepackage{varwidth}
\usepackage{booktabs}
\usepackage{bookmark}

\usepackage{imakeidx}
\makeindex[
  columns=2,
  title=Índice,
  options= -s header.ist,
  intoc
]

\usepackage[
  citestyle=numeric,
  style=numeric,
  backend=biber,
  maxnames=5,
  minnames=3
]{biblatex}

\addbibresource{referências.bib}

\usepackage{hyperref}

\hypersetup{
  pdfinfo={
      Title={Álgebra linear~Verão 2021~Notas das aulas na escola},
      Author={Alex Dantas et al.},
      Keywords={corpos, sistemas lineares},
      Subject={Álgebra linear},
      Producer={TeXstudio 3.0.4 Using Qt Version 5.15.2},
      Creator={LuaHBTeX, Version 1.12.0 (TeX Live 2020/Arch Linux)},
    }
  %	hyperref,
  pdfencoding=auto,
  linktocpage=true,
  colorlinks=true,
  linkcolor=NavyBlue,
  urlcolor=magenta,
  linktoc=all,
  pdfpagelabels,
  bookmarks=true,
  unicode,
  pdfborder = {0 0 0}
  %	filecolor = red,
}

\backgroundsetup{
  scale=1,
  angle=0,
  opacity=.1,
  contents={
      \begin{tikzpicture}[remember picture,overlay]
        \node at ([yshift=11pt,xshift=5pt]current page.center) {\includegraphics[width=8.5cm]{unb}};
      \end{tikzpicture}}
}

\theoremstyle{definition}
\newtheorem{definition}{Definição}[chapter]
\newtheorem{proposition}{Proposição}[chapter]
\newtheorem{example}{Exemplo}[chapter]
\newtheorem{exercise}{Exercício}[chapter]
\newtheorem{remark}{Observação}[chapter]

\newcommand{\MVAt}{{\usefont{U}{mvs}{m}{n}\symbol{`@}}}
\renewcommand{\qedsymbol}{\(\blacksquare\)}
\renewcommand{\solutiontitle}{\noindent\textbf{Solução}\par\noindent}
\definecolor{SolutionColor}{rgb}{0.2,0.9,1}
\colorsolutionboxes
\definecolor{SolutionBoxColor}{rgb}{0.2,0.9,1}
\vqword{Questão}
\vpword{Pontos}
\vsword{Pontuaçõe}

\setkomafont{title}{\bfseries\sffamily\Huge\color{DarkBlue}}
\setkomafont{subtitle}{\bfseries\sffamily\Large\color{DarkBlue}}
\setkomafont{author}{\bfseries\sffamily\large\color{DarkBlue}}
\setkomafont{date}{\bfseries\sffamily\color{DarkBlue}}
\setkomafont{chapter}{\LARGE\color{DarkBlue}}
\setkomafont{section}{\Large\color{DarkBlue}}
\setkomafont{subsection}{\large\color{DarkBlue}}

\urlstyle{rm}

\setmonofont{Consolas}
\newunicodechar{∈}{\makebox[\fontcharwd\font`a]{$\in$}}
\newunicodechar{⊗}{\makebox[\fontcharwd\font`a]{$\otimes$}}

\title{Álgebra Linear II}

\subtitle{
  \href{https://www.mat.unb.br/verao2021/verao/verao_pt.html}{
    XLIX Escola de Verão en Matemática da UnB
  }
}

\titlehead{\Huge
\[
  J_{j}=
  \begin{pmatrix}
    \lambda_{j} & 1           &        &             & 0           \\
                & \lambda_{j} & 1      &             &             \\
                &             & \ddots & \ddots      &             \\
                &             &        & \lambda_{j} & 1           \\
    0           &             &        &             & \lambda_{j}
  \end{pmatrix}
  \in\mathbb{C}^{s_{j}\times s_{j}}.
\]
}

\author{Aulas do professor Alex Carrazedo Dantas\thanks{alex\MVAt mat.unb.br}}

\date{
  Última modificação: \today{} às $\DTMcurrenttime$.
  \\[\baselineskip]
  \url{https://carlosal1015.github.io/Algebra-linear-II/main.pdf}
}

\begin{document}

\maketitle
\tableofcontents

\chapter*{Introdução ao curso\quad$\left(04/01/2021\right)$}

O professor \href{http://buscatextual.cnpq.br/buscatextual/visualizacv.do?id=K4463841D3}{Alex Carrazedo Dantas} é especialista no \href{https://pt.wikipedia.org/wiki/Teoria_dos_grupos}{
  \emph{Teoria dos grupos}}.
Em um curso presencial você pode discutir mais, enquanto em um curso remoto, cada aula tem um pdf \href{https://moodle.mat.unb.br/20201}{
  Moodle MAT} e uma gravação da sessão.
Se você tiver dúvidas sobre o moodle, peça ajuda a
\href{http://buscatextual.cnpq.br/buscatextual/visualizacv.do?id=K8136975Z1}{Carol Lafetá}\footnote{\href{mailto:lafeta.carol@gmail.com}{lafeta.carol\MVAt gmail.com}}.
\section*{Ementa}

\begin{enumerate}\bfseries
  \begin{multicols}{2}
    \item

    Sistemas lineares e matrizes.

    \item

    Espaços vetoriais e transformações lineares.

    \item

    Polinômios e determinantes

    \item

    Decomposicões primárias e formas racionais e de Jordan.

    \item

    Produto interno e teorema espectral.

    \item

    Formas multilineares.
  \end{multicols}
\end{enumerate}

\section*{Critério de avaliação}
%https://introducaocomunicacao.wordpress.com/avaliacao/
\begin{table}[ht!]
  \centering
  \renewcommand\arraystretch{0.2}
  \begin{tabular}{c|c}
    % \toprule
    \textbf{Menção em disciplina} & \textbf{Equivalência numérica} \\
    \midrule
    Superior (SS)                 & $9-10$                         \\
    Média Superior (MS)           & $7-8.9$                        \\
    Média (MM)                    & $5-6.9$
    % \bottomrule
  \end{tabular}
\end{table}
Serão aplicadas $2$ provas, de acordo com o cronograma abaixo, as
quais serão atribuı́das as notas $x$ e $y$.
\[
  \text{MF}=\frac{x+3y}{4}.
\]
O aluno deverá obter média final igual ou superior a $5$ pontos e
$75\%$ de frequência para ser aprovado.
\section*{Tutores}

\begin{itemize}
  \begin{multicols}{3}
    \item\href{https://www.escavador.com/sobre/7335541/sara-raissa-silva-rodrigues}{Sara Raissa Silva Rodrigues}.
    \item \href{https://www.escavador.com/sobre/5504138/geraldo-herbert-beltrao-de-souza}{Geraldo Herbert Beltrão de Souza}.
    \item\href{https://www.escavador.com/sobre/5634752/mattheus-pereira-da-silva-aguiar}{Mattheus Pereira da Silva Aguiar}.
  \end{multicols}
\end{itemize}

\vfill
\nocite{*}
\printbibliography[
  title={\textcolor{DarkBlue}{Referências bibliográficas}{\fontspec[Renderer=Harfbuzz]{NotoColorEmoji.ttf}📚}},
  heading=bibintoc
]

\setpartpreamble{\Huge

\

\[
  J_{j}=
  \begin{pmatrix}
    \lambda_{j} & 1           &        &             & 0           \\
                & \lambda_{j} & 1      &             &             \\
                &             & \ddots & \ddots      &             \\
                &             &        & \lambda_{j} & 1           \\
    0           &             &        &             & \lambda_{j}
  \end{pmatrix}
  \in\mathbb{C}^{s_{j}\times s_{j}}.
\]
}
\part{Teoria}

\chapter{Corpos}

\begin{definition}[Corpo]
	Um \emph{corpo} é um conjunto não vazio $\mathbb{F}$ munido de duas
	operacões: adição mais e multiplicação
	\[ +\colon\mathbb{F}\times\mathbb{F}\left(x,y\right)\mapsto x+y \]
	\[ \cdot\colon\mathbb{F}\times\mathbb{F}\left(x,y\right)\mapsto x\cdot y \]
	e tais que
	\begin{enumerate}
		\item[Asociatividade na adição]

		      $\left(x+y\right)+z=x+\left(y+z\right)$, $\forall x,y,z\mathbb{F}$;

		\item[Existênza de neutro aditivo]

		      $\exists0\in\mathbb{F}$ tal que $x+0=0+x=x$, $\forall x\in\mathbb{F}$;

		\item[Existênza de elemento oposto o inverso aditivo]

		      Dado $x\in\mathbb{F}$, existe $-x\in\mathbb{F}$ tal que $x+\left(-x\right)=\left(-x\right)+x=0$.

		\item[Conmutatividade na adição]

		      $x+y=y+x$, $\forall x,y\in\mathbb{F}$;

		\item[Associatividade na multiplicação]

		      $\left(x\cdot y\right)z=x\cdot\left(y\cdot z\right)$, $\forall x,y,z\in\mathbb{F}$;

		\item[Existênza do elemento neutro na multiplicação]

		      $\exists 1\in\mathbb{F}$ tal que $x\cdot 1=1\cdot x=x$, $\forall x\in\mathbb{F}$;

		\item[Existênza inverso multiplicativo]

		      Dado $x\in\mathbb{F}\setminus\left\{0\right\}$, existe $x^{-1}\in\mathbb{F}$ tal que $x\cdot x^{-1}=x^{-1}\cdot x=1$;

		\item[Conmutatividade na multiplicação]

		      $x\cdot y=y\cdot x$, $\forall x,y\in\mathbb{F}$.

		\item[Distributiva]

		      $x\cdot\left(y+z\right)=xy+xz$, $\forall x,y,z\in\mathbb{F}$.
	\end{enumerate}
\end{definition}

\begin{proposition}
	$x\cot0=0$, $\forall x\in\mathbb{F}$.
\end{proposition}

\begin{example}
	\begin{enumerate}
		\item

		      $\left(\mathbb{Z},+,\cdot\right)$ não é um corpo.

		      O seja el conjunto.

		\item

		      $\left(\mathbb{Q},+,\cdot\right)$ é um corpo, onde
		      $\mathbb{Q}=\left\{\frac{a}{b}\mid a,b\in\mathbb{Z},b\neq0\right\}$.

		\item

		      $\left(\mathbb{C},+,\cdot\right)$ é um corpo, onde
		      $\mathbb{C}=\left\{a+b\imath\mid a,b\in\mathbb{R},\text{ e }\imath^{2}=1\right\}$,
		      $+\colon\left(a+b\imath\right)+\left(c+d\imath\right)=\left(a+c\right)+\left(b+d\right)\imath$
		      $\cdot\colon\left(a+b\imath\right)\cdot\left(c+d\imath\right)=\left(ac-bd\right)+\left(ad+bc\right)\imath$.

		\item

		      $\left(\mathbb{Z}/p\mathbb{Z},+,\cdot\right)$ é um corpo, onde $\mathbb{Z}/p\mathbb{Z}=\left\{\overline{a}\mid \overline{a}\in\mathbb{Z}\right\}$.
	\end{enumerate}
\end{example}
%https://www.learn-portuguese-with-rafa.com/portuguese-keyboard-characters.html


\chapter{Sistemas lineares\quad$\left(07/01/2021\right)$}

\begin{definition}[Sistema linear]
  Um corpo é.
  \[
    \begin{pmatrix}
      a_{11} x_{1}+a_{12} x_{2}+\cdots+a_{1n} x_{n}=y_{2} \\
      a_{21} x_{1}+a_{22} x_{2}+\cdots+a_{2n} x_{n}=y_{2} \\
      \vdots                                              \\
      a_{m1} x_{1}+a_{m2} x_{2}+\cdots+a_{mn} x_{n}=y_{m}
    \end{pmatrix}
  \]
\end{definition}


$c_{1}\left(a_{n}x_{2}+\cdots+a_{1n}x_{n}\right)+
  \cdots+
  c_{m}\left(a_{m}x_{2}+\cdots+a_{mn}x_{n}\right)=
  c_{1}y_{2}+\cdots+c_{m}y_{m}$

$\left(c_{2}a_{11}+\cdots+c_{m}a_{m1}\right)
  x_{1}+\cdots+\left(c_{1}a_{1n}+\cdots+c_{m}a_{w n}\right)
  x_{n}=
  c_{1}y_{2}+\cdots+c_{m}y_{m}$

\[
  \left\{\begin{aligned}
    2 x+3 y-z+w & =5 \\
    x-y+2 z-2 w & =1 \\
    2 x+y+z+w   & =3
  \end{aligned}\right.
\]

\[  \left\{\begin{array}{l}
    x-y+2 z-2 w=1 \\
    2 x+3 y-z+w=5 \\
    2 x+y+z+w=3
  \end{array}\right.
\]

\[
  \left\{\begin{array}{l}
    x-y+2 z-2 w=1 \\
    5 y-5 z+5 w=3 \\
    3 y-3 z+5 w=1
  \end{array}\right.
\]

\[
  \left\{\begin{aligned}
    x-y+2 z-2 w & =1     \\
    y-z+w       & =3 / 5 \\
    y-z+5 / 3 w & =1 / 3
  \end{aligned}\right.
\]

\[
  \left\{\begin{aligned}
    x+z-w        & =8 / 5                                                \\
    y-z+\omega   & =3 / 5                                                \\
    2 / 3 \omega & =\frac{1}{3}-\frac{3}{5}=\frac{5-9}{15}=-\frac{4}{15} \\
    \omega       & =-\frac{12}{30}=-\frac{4}{10}=-\frac{2}{5}
  \end{aligned}\right.
\]

\[
  \begin{array}{l}
    \left\{(x, y, z, \omega) \in \mathbb{Q}^{4} \mid x=-z+\frac{6}{5}, y=z+1, z \in \mathbb{Q}, \omega=-\frac{2}{5}\right\}= \\
    =\left\{\left(-z+\frac{6}{5}, z+1, z,-\frac{2}{5}\right) \mid z \in \mathbb{Q}\right\}
  \end{array}
\]

\[
  \left\{\begin{array}{l}
    x+y=2 \\
    x-y=0
  \end{array} \quad\left\{\begin{array}{l}
    x+2 y=5 \\
    x-y=-1
  \end{array}\right.\right.
\]

\[
  f\colon
  \left\{1,\ldots,m\right\}\times
  \left\{1,\ldots,n\right\}\rightarrow
  F=f\left(i,j\right)=a_{ij}.
\]

\[
  \begin{pmatrix}
    f\left(1,1\right) & f\left(1,2\right) & \cdots & f\left(1,n\right) \\
    f\left(2,1\right) & f\left(2,2\right) & \cdots & f\left(2,n\right) \\
    \vdots            &                   & \ddots & \vdots            \\
    f\left(x,1\right) & f\left(w,2\right) & \cdots & f\left(w,n\right)
  \end{pmatrix}\quad
  \begin{pmatrix}
    a_{n}   & a_{12}  & \cdots & a_{1 n} \\
    a_{21}  & a_{22}  & \cdots & a_{2 n} \\
    \vdots  &         & \ddots & \vdots  \\
    a_{n 1} & a_{m 2} & \cdots & a_{w n}
  \end{pmatrix}
\]
\chapter{Matrizes\quad$\left(08/01/2021\right)$}

Podemos denotar uma matriz $A$ sobre um corpo $\mathbb{F}$ de ordem
$m\times n$ por $A={\left(a_{ij}\right)}_{m\times n}$.

Sejam $A={\left(a_{ij}\right)}_{m\times n}$ e
$B={\left(b_{jl}\right)}_{n\times p}$ duas matrizes sobre um corpo
$\mathbb{F}$.
Definimos o producto de $A$ por $B$ como a matriz
$C={\left(c_{il}\right)}_{m\times p}$ dada por
\[
  c_{il}=
  \sum_{j=1}^{n}
  a_{i j}b_{jl}=
  a_{i1}b_{il}+\cdots+a_{in}b_{nl}
\]

\section*{Ilustração}

\[
  \begin{pmatrix}
    a_{11} & \cdots & a_{1n} \\
    \vdots & \vdots & \vdots \\
    a_{m1} & \cdots & a_{mn}
  \end{pmatrix}
  \begin{pmatrix}
    b_{11} & \cdots & b_{1p} \\
    \vdots & \vdots & \vdots \\
    b_{n1} & \cdots & b_{np}
  \end{pmatrix}=
  \begin{pmatrix}
    c_{11} & \cdots & c_{1p} \\
    \vdots & \vdots & \vdots \\
    c_{m1} & \cdots & c_{mp}
  \end{pmatrix}.
\]

\begin{example}
  Temos que
  \[
    {\begin{pmatrix}
          1 & 2  \\
          2 & -1
        \end{pmatrix}}_{2\times2}
      {\begin{pmatrix}
          1 & 0 & 4 \\
          3 & 2 & 5
        \end{pmatrix}}_{2\times3}=
    {\begin{pmatrix}
      7  & 4  & 14 \\
      -1 & -2 & 3
    \end{pmatrix}}_{2\times3}.
  \]
\end{example}

\begin{proposition}
  Sejam matrices $A={\left(a_{ij}\right)}_{m\times n}$,
  $B={\left(a_{jl}\right)}_{n\times p}$ e
  $C={\left(a_{lk}\right)}_{p\times q}$ matrizes sobre um corpo
  $\mathbb{F}$.
  Então $\left(AB\right)C=A\left(BC\right)$.
\end{proposition}

\begin{proof}
  Veja que
  \begin{math}
    \left(AB\right)C=
    {\left(\alpha_{ik}\right)}_{m\times q}
  \end{math},
  \begin{math}
    AB=
    {\left(d_{il}\right)}_{m\times p}
  \end{math}
  onde
  \begin{align*}
    d_{il}
     & =\sum_{l=1}^{n}
    a_{ij}b_{jl}
    \shortintertext{e}
    \alpha_{ik}
     & =\sum_{l=1}^{p}
    d_{il}c_{lk}=
    \sum_{l=1}^{p}
    \left(\sum_{j=1}^{n}a_{ij}b_{jl}\right)c_{lk}= \\
     & =\sum_{l=1}^{p}
    \left(\sum_{j=1}^{n}a_{ij}b_{jl}c_{lk}\right)= \\
     & =\sum_{j=1}^{n}
    a_{ij}\left(\sum_{l=1}^{p}b_{jl}c_{lk}\right)=\beta_{ik}
  \end{align*}
  com $A\left(BC\right)={\left(\beta_{ik}\right)}_{m\times q}$.
\end{proof}

Chamaremos a matriz quadrada
$I_{m}={\left(\delta_{ij}\right)}_{m\times m}$ definida por
\[
  \delta_{ij}=
  \begin{cases}
    1 & , \text{se }i=j,     \\
    0 & , \text{se }i\neq j,
  \end{cases}
\]
de matriz identidade de ordem $m\times m$.

Note que se $A={\left(a_{jl}\right)}_{m\times n}$, então $I_{m}A=A$, e se $B={\left(b_{li}\right)}_{n\times m}$, então $BI_{m}=B$.


$I_{m}A={\left(c_{il}\right)}_{m\times m}$ é tal que
$c_{il}=\sum_{j=1}^{m}\delta_{ij}a_{jl}=a_{il}$ con $1\leq i\leq m$,
e $I_{m}A=A$.

\begin{example}
  Se $m=3$, então
  \[
    I_{3}=
    \begin{pmatrix}
      1 & 0 & 0 \\
      0 & 1 & 0 \\
      0 & 0 & 1
    \end{pmatrix}.
  \]
\end{example}
Dizemos que uma matriz quadrada $A={\left(a_{ij}\right)}_{m\times m}$
tem inversa se existe uma matriz
$B={\left(b_{ij}\right)}_{m\times m}$ tal que $AB=BA=I_{m}$.

Denotaremos a matriz $B$ por $A^{-1}$.

\begin{definition}
  Seja $c\in\mathbb{F}\setminus\left\{0\right\}$.
  Uma matriz quadrada de ordem $m\times m$ $E$ é dita elementar se
  $E$ é de uma das formas
  \begin{enumerate}
    \item

          $E_{1}={\left(e_{ij}\right)}_{m\times m}$, onde
          \[
            e_{ij}=
            \begin{cases}
              \delta_{ij}, & \text{se }i\neq k \\
              \delta_{ij}, & \text{se }i=k
            \end{cases}\qquad
            \fbox{\begin{varwidth}{\textwidth}
                $m=3$, $k=2$\\[\baselineskip]
                \begin{math}
                  E_{1}=
                  \begin{pmatrix}
                    1 & 0 & 0 \\
                    0 & c & 0 \\
                    0 & 0 & 1
                  \end{pmatrix}
                \end{math}
              \end{varwidth}}
          \]
          con $k$ um inteiro fixo entre $1$ e $m$;
    \item

          $E_{2}={\left(e_{ij}\right)}_{m\times m}$, onde
          \[
            e_{ij}=
            \begin{cases}
              \delta_{ij}, & \text{se }i\neq l\text{ e }i\neq k \\
              \delta_{lj}, & \text{se }i=k                      \\
              \delta_{kj}, & \text{se }i=l
            \end{cases}\qquad
            \fbox{\begin{varwidth}{\textwidth}
                $m=3$, $k=2$, $l=3$\\[\baselineskip]
                \begin{math}
                  \begin{pmatrix}
                    1 & 0 & 0 \\
                    0 & 0 & 1 \\
                    0 & 1 & 0
                  \end{pmatrix}
                \end{math}
              \end{varwidth}}
          \]
          con $k<l$ inteiros fixos entre $1$ e $m$;
    \item

          $E_{3}={\left(e_{ij}\right)}_{m\times m}$, onde
          \[
            e_{ij}=
            \begin{cases}
              \delta_{i j},                  & i\neq k \\
              \delta_{kj}+c\cdot\delta_{lj}, & i=k
            \end{cases}\qquad
            \fbox{\begin{varwidth}{\textwidth}
                $m=3$, $k=2$, $l=3$\\[\baselineskip]
                \begin{math}
                  \begin{pmatrix}
                    1 & 0 & 0 \\
                    0 & 1 & c \\
                    0 & 0 & 1
                  \end{pmatrix}
                \end{math}
              \end{varwidth}}
          \]
  \end{enumerate}
\end{definition}

\begin{example}
  Calcule
  \[
    \begin{pmatrix}
      1 & 0 & 0 \\
      0 & 1 & 2 \\
      0 & 0 & 1
    \end{pmatrix}
    \begin{pmatrix}
      1 & 2 & 3  & -1 \\
      2 & 2 & 1  & 1  \\
      1 & 1 & -1 & 2  \\
    \end{pmatrix}=
    \begin{pmatrix}
      1 & 2 & 3  & -1
      4 & 4 & -1 & 5
      1 & 1 & -1 & 2
    \end{pmatrix}
  \]
\end{example}
Dada uma matriz $A={\left(a_{ij}\right)}_{m\times m}$ o efeito de multiplicar uma matriz elementar $E$ por $A$ pode ser colocado como:
\begin{enumerate}
  \item

        $E_{1}A$: multiplica uma linha $k$ de $A$ por um escalar $c$;

  \item

        $E_{2}A$: troca duas linhas $l$ e $k$ de posições
        $\left(k<l\right)$;

  \item

        $E_{3}A$: soma uma linha $k$ com outra linha $l$ multiplicada por um escalar $c\in\mathbb{F}$.
\end{enumerate}
\[
  \begin{array}{l}
    E_{1}=\left(e_{i j}\right)_{m \times m} \equiv A=\left(a_{j \ell}\right)_{m \times n} \\
    E_{1} A=\left(c_{i \ell}\right)_{m \times n} \text { com }
  \end{array}
\]
\chapter{Aula de reposição\quad$\left(09/01/2021\right)$}

\begin{definition}[Matriz reducida por linhas]
  Uma matriz $A={\left(a_{ij}\right)}_{m\times n}$ sobre $\mathbb{F}$
  é deja reduzida por linhas se

  \begin{enumerate}
    \item

          O primeiro elemento não nulo de cada linha não nula é
          igual $1$;

    \item

          cada columna que possui o primeiro elemento não nulo de

    \item

          uma linha não possui todos os outros elementos iguais a
          $0$;
  \end{enumerate}


  Sea além disso, esa matriz $A$ satisfaz

  \begin{enumerate}

    \item

          Todas linhas nulas ocorrem abaixo das linhas não nulas;

    \item

          Se $1,\ldots,r$ $\left(r\leq m\right)$ são as linhas não
          nulas de $A$ com os primeiros elementos não nunos ocurrendo
          nas colunas $k_{1},k_{2},k_{r}$, respectivamente, então
          $k_{1}<k_{2}<\cdots<k_{r}$, dizemos que $A$ está na forma
          escada reduzida.
  \end{enumerate}
\end{definition}

\begin{example}
  \begin{enumerate}
    \item As seguintes matrizes estão na forma reduzida:
          \begin{enumerate}
            \item
                  \begin{math}
                    \begin{pmatrix}
                      0 & 1 & 2 & 3 \\
                      1 & 0 & 1 & 1
                    \end{pmatrix},
                  \end{math}.
            \item

                  \begin{math}
                    \begin{pmatrix}
                      0 & 0 & 0 \\
                      1 & 0 & 1 \\
                      0 & 1 & 0
                    \end{pmatrix},
                  \end{math}.

            \item

                  \begin{math}
                    \begin{pmatrix}
                      0 & 0 & 0 \\
                      0 & 0 & 0 \\
                      0 & 0 & 0
                    \end{pmatrix},
                  \end{math}.

            \item

                  \begin{math}
                    \begin{pmatrix}
                      1 & 0 & 0 & 0 \\
                      0 & 1 & 1 & 2 \\
                      0 & 0 & 0 & 0
                    \end{pmatrix},
                  \end{math}.
          \end{enumerate}

    \item As seguientes matrizes estão na forma escada reducida

          \begin{enumerate}
            \item

                  \begin{math}
                    \begin{pmatrix}
                      1 & 0 & 2 & 0 \\
                      0 & 1 & 3 & 0 \\
                      0 & 0 & 0 & 1
                    \end{pmatrix}
                  \end{math}.

            \item

                  \begin{math}
                    \begin{pmatrix}
                      1 & 0 \\
                      0 & 1 \\
                      0 & 0
                    \end{pmatrix}
                  \end{math}.


            \item

                  \begin{math}
                    \begin{pmatrix}
                      1 & 2 & 3 & 0 & 1 \\
                      0 & 0 & 0 & 1 & 0
                    \end{pmatrix}
                  \end{math}.
          \end{enumerate}

  \end{enumerate}
\end{example}

\begin{remark}
  .
\end{remark}
\input{aulas/espaços-vetoriais}

\setpartpreamble{%
  \begin{center}
    \includesvg[width=.4\paperwidth]{The_four_subspaces-X-1}
  \end{center}
}
\part{Prática}

\chapter{Exercícios de Fixação\quad$\left(08/01/2021\right)$}

\begin{questions}
  \question\label{exercício:1.1}

  Seja $\mathbb{F}$ um corpo.
  Dizemos que um subconjunto $\mathbb{K}$ de $\mathbb{F}$ é um
  subcorpo de $\mathbb{F}$ se $\mathbb{K}$ munido das operações de
  adição e multiplicação de $\mathbb{F}$ é um corpo.
  Mostre que os seguintes subconjuntos são subcorpos de $\mathbb{C}$.

  \begin{parts}
    \begin{multicols}{3}
      \part\label{exercício:1.1a}

      \begin{math}
        \mathbb{Q}
        \left(\sqrt{3}\right)=
        \left\{
        a+b\sqrt{3}\mid a,b\in\mathbb{Q}
        \right\}
      \end{math};

      \part\label{exercício:1.1b}

      \begin{math}
        \mathbb{Q}
        \left(i\right)=
        \left\{
        a+bi\mid a,b\in\mathbb{Q}
        \text{ e }i^{2}=-1
        \right\}
      \end{math};

      \part\label{exercício:1.1c}

      \begin{math}
        \mathbb{Q}
        \left(i\sqrt{2}\right)=
        \left\{
        a+bi\sqrt{2}\mid a,b\in\mathbb{Q}
        \text{ e }i^{2}=-1
        \right\}
      \end{math}.
    \end{multicols}
  \end{parts}

  \begin{solutionordottedlines}
    \begin{parts}
      \part
      .

      \part

      .

      \part

      .
    \end{parts}
  \end{solutionordottedlines}

  \question\label{exercício:1.2}

  Mostre que:

  \begin{parts}
    \part\label{exercício:1.2a}

    Todo subcorpo de $\mathbb{C}$ tem $\mathbb{Q}$ como subcorpo;

    \part\label{exercício:1.2b}

    Todo corpo de característica $0$ tem uma cópia de $\mathbb{Q}$;

    \part\label{exercício:1.2c}

    Se $\mathbb{K}$ contém propriamente $\mathbb{R}$ e é um subcorpo de
    $\mathbb{C}$, então $\mathbb{K}=\mathbb{C}$.

  \end{parts}

  \begin{solutionordottedlines}
    \begin{parts}
      \part
      .

      \part

      .

      \part

      .
    \end{parts}
  \end{solutionordottedlines}

  \question\label{exercício:1.3}

  Considere o corpo finito com $5$ elementos
  \begin{math}
    \mathbb{Z}/5\mathbb{Z}=
    \left\{
    \overline{0},
    \overline{1},
    \overline{2},
    \overline{3},
    \overline{4}
    \right\}
  \end{math}.

  \begin{parts}
    \part\label{exercício:1.3a}

    Mostre que
    \[
      \mathbb{F}=
      \left\{
      a+bi\mid a,b\in\mathbb{Z}/5\mathbb{Z}
      \text { e }
      i^{2}=
      \overline{3}
      \right\}
    \]

    munido das operações

    \[
      \begin{aligned}
        \overline{+}\colon\mathbb{F}\times\mathbb{F}     & \longrightarrow\mathbb{F}     \\
        \left(\left(a+bi\right),\left(c+di\right)\right) & \longmapsto \left(a+c\right)+
        \left(b+d\right)i
      \end{aligned}\quad
      \begin{aligned}
        \overline{\cdot}\colon\mathbb{F}\times\mathbb{F}  & \longrightarrow\mathbb{F}                   \\
        \left(\left(a+b i\right),\left(c+di\right)\right) & \longmapsto \left(ac+\overline{3}bd\right)+
        \left(ad+bc\right)i
      \end{aligned}
    \]

    é um corpo com 25 elementos;

    \part\label{exercício:1.3b}

    Mostre que $\mathbb{Z}/5\mathbb{Z}$ é um subcorpo de $\mathbb{F}$.
    Qual é a característica de $F$?
  \end{parts}

  \begin{solutionordottedlines}
    \begin{parts}
      \part

      .

      \part

      .
    \end{parts}
  \end{solutionordottedlines}

  \question\label{exercício:1.4}

  Determine o conjunto solução de cada sistema linear dado.

  \begin{parts}%\setlength{\columnsep}{-20pt}
    \begin{multicols}{2}
      \part\label{exercício:1.4a}

      \begin{math}
        \csysteme[xyzw]{
          x-2y+z+w=1,
          2x+y-z=3,
          2x+y-5z+w=4
        }
      \end{math}
      em $\mathbb{R}$,

      \part\label{exercício:1.4b}

      \begin{math}
        \csysteme[xyzw]{
          x-\sqrt{3}y+z+w=1+\sqrt{3},
          \left(2\+\sqrt{3}\right)x+y-z=3,
          2x+y-\left(1\-\sqrt{3}\right)z+w=4
        }
      \end{math}
      em $\mathbb{Q}\left(\sqrt{3}\right)$,

      \part\label{exercício:1.4c}

      \begin{math}
        \csysteme[xyzw]{
          x-2iy+2z-w=0,
          \left(2\+i\right)x+z+w=0,
          2ix+y-5z+\left(1\+i\right)w=0
        }
      \end{math}
      em $\mathbb{C}$,

      \part\label{exercício:1.4d}

      \begin{math}
        \csysteme[xyzw]{
          x-\overline{2}y+\overline{2}z-w=\overline{0},
          \overline{2}x+z+w=\overline{0},
          \overline{2}x+y-\overline{3}z+w=\overline{0}
        }
      \end{math}
      em $\mathbb{Z}/5\mathbb{Z}$,

      \part\label{exercício:1.4e}

      \begin{math}
        \csysteme[xyzw]{
          x-\overline{2}iy+\overline{2}z-w=\overline{0},
          \left(\overline{2}\+i\right)x+z+w=\overline{0},
          \overline{2}ix+y-\overline{3}z+\left(\overline{1}\+i\right)w=\overline{0}
        }
      \end{math}
      em $\mathbb{F}$ de (a) da questão 3.
    \end{multicols}
  \end{parts}%\setlength{\columnsep}{10pt}

  \begin{solutionordottedlines}
    \begin{parts}
      \part

      .

      \part

      .

      \part

      .

      \part

      .

      \part

      .
    \end{parts}
  \end{solutionordottedlines}

  \question\label{exercício:1.5}

  Mostre que se dois sistemas lineares $2\times2$ possuem o mesmo
  conjunto solução, então eles são equivalentes.
  Determine, se existir, dois sistemas lineares $2\times3$ com mesmo
  conjunto solução mas não equivalentes.

  \begin{solutionordottedlines}
  \end{solutionordottedlines}

  \question\label{exercício:1.6}

  Considere o sistema linear sobre $\mathbb{Q}$

  \[
    \csysteme[xyzw]{
      x-2y+z+2w=1,
      x+y-z+w=2,
      x+7y-5z-w=3
    }
  \]

  Mostre que esse sistema não tem solução.

  \begin{solutionordottedlines}
  \end{solutionordottedlines}

  \question\label{exercício:1.7}

  Determine todos $a,b,c,d\in\mathbb{R}$ tais que o sistema linear

  \[
    \begin{bNiceMatrix}
      3  & -6 & 2 & 1 \\
      -2 & 4  & 1 & 3 \\
      0  & 0  & 1 & 1 \\
      1  & -2 & 1 & 0
    \end{bNiceMatrix}
    \cdot
    \begin{bNiceMatrix}
      x \\
      y \\
      z \\
      w
    \end{bNiceMatrix}=
    \begin{bNiceMatrix}
      a \\
      b \\
      c \\
      d
    \end{bNiceMatrix}
  \]
  tem solução.

  \begin{solutionordottedlines}
  \end{solutionordottedlines}

  \question\label{exercício:1.8}

  Encontre duas matrizes $A$ e $B$ de ordens iguais a $3\times3$ tais
  que $AB$ é uma matriz nula mas $BA$ não é.

  \begin{solutionordottedlines}
  \end{solutionordottedlines}

  \question\label{exercício:1.9}

  Mostre que toda matriz elementar é inversível e calcule a inversa
  de cada tipo.

  \begin{solutionordottedlines}
  \end{solutionordottedlines}

  \question\label{exercício:1.10}

  Determine a matriz inversa da matriz

  \[
    A=
    \begin{bNiceMatrix}
      1 & 2 & 3 & 4 \\
      0 & 2 & 3 & 4 \\
      0 & 0 & 3 & 4 \\
      0 & 0 & 0 & 4
    \end{bNiceMatrix}.
  \]

  \begin{solutionordottedlines}
  \end{solutionordottedlines}


  \question\label{exercício:11.1}

  Considere a matriz

  \NiceMatrixOptions{
    cell-space-top-limit = 3pt,
    %cell-space-bottom-limit = 2pt
  }

  \[
    A=
    \begin{bNiceMatrix}
      \overline{1} & \overline{2} & \overline{3} & \overline{4} \\
      \overline{0} & \overline{2} & \overline{3} & \overline{4} \\
      \overline{0} & \overline{0} & \overline{3} & \overline{4} \\
      \overline{0} & \overline{0} & \overline{0} & \overline{4}
    \end{bNiceMatrix}
  \]

  \NiceMatrixOptions{
    cell-space-top-limit = 0pt,
  }

  com entradas no corpo com cinco elementos
  \begin{math}
    \mathbb{Z}/5\mathbb{Z}=
    \left\{
    \overline{0},
    \overline{1},
    \overline{2},
    \overline{3},
    \overline{4}
    \right\}
  \end{math}.
  Calcule sua inversa.

  \begin{solutionordottedlines}
  \end{solutionordottedlines}

  \question\label{exercício:1.12}

  Considere a matriz
  \[
    A=
    \begin{bNiceMatrix}
      1  & 2  & 1 & 0 \\
      -1 & 0  & 3 & 5 \\
      1  & -2 & 1 & 1
    \end{bNiceMatrix}.
  \]
  Encontre uma matriz na forma e uma matriz invertível $P$ tal que
  $R=PA$.

  \begin{solutionordottedlines}
  \end{solutionordottedlines}
\end{questions}

\chapter{Exercícios de Fixação\quad$\left(15/01/2021\right)$}

\begin{questions}
	\question\label{exercício:2.1}

	Defina sobre $\mathbb{R}^{2}$ as seguintes operações:

	\[
		\begin{aligned}
			+\colon\mathbb{R}^{2}\times\mathbb{R}^{2}      & \longrightarrow\mathbb{R}^{2} \\
			\left(\left(x,y\right),\left(a,b\right)\right) & \longmapsto
			\left(x+a,0\right)
		\end{aligned}\qquad
		\begin{aligned}
			\cdot\colon\mathbb{R}\times\mathbb{R}^{2} & \longrightarrow\mathbb{R}^{2} \\
			\left(c,\left(x,y\right)\right)           & \longmapsto
			\left(cx,0\right)
		\end{aligned}
	\]

	O conjunto $\mathbb{R}^{2}$ é um espaço vetorial com essas operações?

	\begin{solutionordottedlines}
	\end{solutionordottedlines}

	\question\label{exercício:2.2}

	Defina sobre
	\begin{math}
		V=
		\left\{
		\left(x,y\right)\in
		\mathbb{R}^{2}\mid
		x>0,y>0
		\right\}
	\end{math}
	as seguintes operações:

	\[
		\begin{aligned}
			+\colon\mathbb{R}^{2}\times\mathbb{R}^{2}      & \longrightarrow\mathbb{R}^{2} \\
			\left(\left(x,y\right),\left(a,b\right)\right) & \longmapsto
			\left(xa,yb\right)
		\end{aligned}\qquad
		\begin{aligned}
			\cdot\colon\mathbb{R}\times\mathbb{R}^{2} & \longrightarrow\mathbb{R}^{2} \\
			\left(c,\left(x,y\right)\right)           & \longmapsto
			\left(x^{c},y^{c}\right)
		\end{aligned}
	\]
	Mostre que $V$ é um espaço vetorial com essas operações.

	\begin{solutionordottedlines}
	\end{solutionordottedlines}

	\question\label{exercício:2.3}

	Resolva:
	\begin{parts}
		\part\label{exercício:2.3a}

		O vetor $\left(3,-1,0,-1\right)$ pertence ao subespaço
		\begin{math}
			W=
			\left\langle
			\left(2,-1,3,2\right),
			\left(-1,1,1,3\right),
			\left(1,1,9,-5\right)
			\right\rangle
		\end{math}
		de $\mathbb{R}^{4}$?

		\part\label{exercício:2.3b}

		Determine uma base para o subespaço vertorial de $\mathbb{R}^{5}$
		das soluções do sistema linear homogêneo
		$\csysteme[xyzwt]{x-2y+z+w+t=0,3x+y-z-4t=0,2x+y-3z+w=0}$;

		\part\label{exercício:2.3c}

		Determine uma base para o subespaço vertorial de
		${\left(\mathbb{Z}/5\mathbb{Z}\right)}^{5}$ das soluções do
		sistema linear homogêneo
		\begin{math}
			\csysteme[xyzwt]{
				x-\overline{2}y+z+w+t=\overline{0},
				\overline{2}x+y-z+t=\overline{0},
				\overline{3}x+y+\overline{3}z+w=\overline{0}
			}
		\end{math},
	\end{parts}

	\begin{solutionordottedlines}
		\begin{parts}
			\part

			.

			\part

			.

			\part

			.
		\end{parts}
	\end{solutionordottedlines}

	\question\label{exercício:2.4}

	Sejam $V$ um espaço vetorial sobre um corpo $F$ e $U$ e $W$
	subespaços de $V$ tais que $U+W=V$ e $U \cap W=\left\{0\right\}$.
	Mostre que cada vetor $v\in V$ é escrito de maneira única como
	$v=u+w$, onde $u\in U$ e $w\in W$.

	\begin{solutionordottedlines}
	\end{solutionordottedlines}

	\question\label{exercício:2.5}

	Mostre que o conjunto dos polinômios sobre uma variável com
	coeficientes em $\mathbb{R}$ é um espaço vetorial sobre
	$\mathbb{R}$ munido das operações usuais de soma e multiplicação
	por escalar.
	Determine uma base para esse espaço vetorial.

	\begin{solutionordottedlines}
	\end{solutionordottedlines}

	\question\label{exercício:2.6}

	Seja $S$ um subconjunto de um espaço vetorial $V$.
	Mostre que $S$ é LD se, e somente se, existir um vetor $v\in S$ que
	pode ser escrito como combinação linear dos elementos de
	$S\setminus\left\{v\right\}$.

	\begin{solutionordottedlines}
	\end{solutionordottedlines}

	\question\label{exercício:2.7}

	Considere o seguinte espaço vetorial sobre $\mathbb{R}$:

	\[
		\mathcal{P}_{3}
		\left(\mathbb{R}\right)=
		\left\{
		a+bx+cx^{2}+dx^{3}\mid
		a,b,c,d\in\mathbb{R}
		\right\}.
	\]

	\begin{parts}
		\part\label{exercício:2.7a}

		Mostre que $\alpha=\left\{1,2+x,3x-x^{2},x-x^{3}\right\}$ é uma
		base de $\mathcal{P}_{3}\left(\mathbb{R}\right)$;

		\part\label{exercício:2.7b}

		Escreva as coordenadas de $p\left(x\right)=1+x+x^{2}+x^{3}$ com
		relação a base $\alpha$;

		\part\label{exercício:2.7c}

		Determine as matrizes mudança de base
		${\left[I\right]}_{\alpha}^{e}$ e
		${\left[I\right]}_{e}^{\alpha}$, onde
		$e=\left\{1,x,x^{2},x^{3}\right\}$.
	\end{parts}

	\begin{solutionordottedlines}
		\begin{parts}
			\part

			.

			\part

			.

			\part

			.
		\end{parts}
	\end{solutionordottedlines}

	\question\label{exercício:2.8}

	Faça o que se pede:

	\begin{parts}
		\part\label{exercício:2.8a}

		Considere a função
		$T\colon\mathbb{C}\to\mathcal{M}_{2}\left(\mathbb{R}\right)$
		dada por

		\[
			T\left(x+yi\right)=
			\begin{bmatrix*}
				x+7y & 5y \\
				-10y & x-7y
			\end{bmatrix*}.
		\]

		Moste que $T$ é uma transformação linear.
		Prove que
		\begin{math}
			T\left(z_{1}z_{2}\right)=
			T\left(z_{1}\right)
			T\left(z_{2}\right),
			\forall z_{1},z_{2}\in\mathbb{C}
		\end{math};

		\part\label{exercício:2.8b}

		Mostre que a composta de transformações lineares é uma transformação linear;

		\part\label{exercício:2.8c}

		Mostre que uma função $T\colon\mathbb{F}^{n}\to\mathbb{F}$ é uma
		transformação linear se, e somente se, existem escalares
		$c_{1},\dotsc,c_{n}$ no corpo $\mathbb{F}$ tais que

		\[
			T
			\left(x_{1},\dotsc,x_{n}\right)=
			c_{1}x_{1}+\dotsc+c_{n}x_{n}.
		\]
	\end{parts}

	\begin{solutionordottedlines}
		\begin{parts}
			\part

			.

			\part

			.

			\part

			.
		\end{parts}
	\end{solutionordottedlines}

	\question\label{exercício:2.9}

	Faça o que se pede:

	\begin{parts}
		\part\label{exercício:2.9a}

		Considere $\mathbb{R}^{4}$ e seus subespaços
		\begin{math}
			W=
			\left\langle
			\left(1,0,1,1\right),
			\left(0,-1,-1,-1\right)
			\right\rangle
		\end{math}
		e
		\begin{math}
			U=
			\left\{
			\left(x,y,z,w\right)\in
			\mathbb{R}^{4}\mid
			x+y=0,
			z+t=0
			\right\}
		\end{math}.
		Determine uma transformação linear
		$T\colon\mathbb{R}^{4}\to\mathbb{R}^{4}$ tal que
		$\operatorname{Nuc}\left(T\right)=U$ e
		$\operatorname{Im}\left(T\right)=W$;

		\part\label{exercício:2.9b}

		Considere ${\left(\mathbb{Z}/5\mathbb{Z}\right)}^{4}$ e seus
		subespaços
		\begin{math}
			W=
			\left\langle
			\left(
			\overline{1},
			\overline{0},
			\overline{1},
			\overline{1}
			\right),
			\left(
			\overline{0},
			\overline{4},
			\overline{4},
			\overline{4}
			\right)
			\right\rangle
		\end{math}
		e
		\begin{math}
			U=
			\left\{
			\left(x,y,z,w\right)\in
			{\left(\mathbb{Z}/5\mathbb{Z}\right)}^{4}\mid
			x+y=
			\overline{0},
			z+t=
			\overline{0}
			\right\}
		\end{math}.
		Determine uma transformação linear
		$T\colon\mathbb{R}^{4}\to\mathbb{R}^{4}$ tal que
		$\operatorname{Nuc}\left(T\right)=V$ e
		$\operatorname{Im}\left(T\right)=W$;

		\part\label{exercício:2.9c}

		Determine uma base para o núcleo e uma base para a imagem da
		transformação linear $T\colon\mathbb{C}^{2}\to\mathbb{R}^{2}$
		dada por $T\left(x+yi,a+bi\right)=\left(x+2a,-x+2b\right)$.
	\end{parts}

	\begin{solutionordottedlines}
		\begin{parts}
			\part

			.

			\part

			.
			\part

			.
		\end{parts}
	\end{solutionordottedlines}

	\question\label{exercício:2.10}

	Sejam $V$ e $W$ espaços vetoriais sobre um mesmo corpo $\mathbb{F}$
	e $T\colon V\to W$.
	Mostre que $T$ é injetora se, e somente se, $T$ leva subconjunto LI
	em subconjunto LI.

	\begin{solutionordottedlines}
	\end{solutionordottedlines}

	\question\label{exercício:2.11}

	Seja
	$T\colon\mathbb{C}^{3}\to\mathcal{P}_{2}\left(\mathbb{C}\right)$
	a transformação linear definida por
	$T\left(1,0,0\right)=1+ix^{2}$,
	$T\left(0,1,0\right)=x+x^{2}$ e
	$T\left(0,0,1\right)=i+x$.
	Exiba uma fórmula para $T$ e decida se $T$ é um isomorfismo.

	\begin{solutionordottedlines}
	\end{solutionordottedlines}

	\question\label{exercício:2.12}

	Seja $F$ um corpo e $T\colon\mathbb{F}^{2}\to\mathbb{F}^{2}$ dada
	por
	\begin{math}
		T\left(x,y\right)=
		\left(x+y,x\right),
		\forall\left(x,y\right)\in
		\mathbb{F}^{2}
	\end{math}.
	Mostre que $T$ é um isomorfismo e exiba uma fómula para $T^{-1}$.

	\begin{solutionordottedlines}
	\end{solutionordottedlines}

	\question\label{exercício:2.13}

	Considere as bases $\alpha=\left\{1,1+x,1+x^{2}\right\}$ de
	$\mathcal{P}_{2}\left(\mathbb{R}\right)$ e
	\begin{math}
		\beta=
		\left\{
		\left(1,0\right),
		\left(i,0\right),
		\left(1,1\right),
		\left(1,i\right)
		\right\}
	\end{math}
	de $\mathbb{C}^{2}$ como espaços vetoriais sobre $\mathbb{R}$.
	Determine as coordenadas da transformação linear
	$T\colon\mathcal{P}_{2}\left(\mathbb{R}\right)\to\mathbb{C}^{2}$ dada
	por $T\left(a+bx+cx^{2}\right)=\left(a+bi,b+ci\right)$ com relação
	à base de
	\begin{math}
		L
		\left(
		\mathcal{P}_{2}\left(\mathbb{R}\right),
		\mathbb{C}^{2}
		\right)
	\end{math}
	construída no Teorema 2-(ii) da Aula $9$.

	\begin{solutionordottedlines}
	\end{solutionordottedlines}

	\question\label{exercício:2.14}

	Considere a base
	\begin{math}
		\alpha=
		\left\{
		\left(1,0,-1\right),
		\left(1,1,1\right),
		\left(2,2,0\right)
		\right\}
	\end{math}
	de $\mathbb{C}^{3}$ como espaço vetorial sobre $\mathbb{C}$.
	Determine a base dual $\alpha^{\ast}$ de
	${\left(\mathbb{C}^{3}\right)}^{\ast}$.

	\begin{solutionordottedlines}
	\end{solutionordottedlines}

	\question\label{exercício:2.15}

	Considere $T\colon\mathbb{R}^{2}\to\mathbb{R}^{3}$ dada por
	$T\left(x,y\right)=\left(2x+3y,y-x,3x\right)$ e as bases
	$\alpha=\left\{\left(1,2\right),\left(2,-1\right)\right\}$ de
	$\mathbb{R}^{2}$ e
	\begin{math}
		\beta=
		\left\{
		\left(1,1,1\right),
		\left(0,1,1\right),
		\left(0,0,1\right)
		\right\}
	\end{math}
	de $\mathbb{R}^{3}$.

	Calcule ${\left[T\right]}_{\beta}^{\alpha}$,
	${\left[T\right]}_{\beta}^{e_{1}}$ e
	${\left[T\right]}_{e_{2}}^{\alpha}$ onde $e_{1}$ é a base canônica
	de $\mathbb{R}^{2}$ e $e_{2}$ é a base canônica de
	$\mathbb{R}^{3}$.

	\begin{solutionordottedlines}
	\end{solutionordottedlines}

	\question\label{exercício:2.16}

	Sejam
	$T\colon\mathbb{R}^{3}\to\mathcal{P}_{2}\left(\mathbb{R}\right)$ e
	$G\colon\mathcal{P}_{2}\left(\mathbb{R}\right)\to\mathbb{R}^{3}$
	transformações lineares tais que
	\[
		{\left[T\right]}_{\beta}^{\alpha}=
		\begin{bmatrix*}
			1 & 2 & -1 \\
			1 & 0 & -1 \\
			0 & -1 & 0
		\end{bmatrix*}\quad
		\text { e }\quad
		{\left[G\right]}_{\alpha}^{\beta}=
		\begin{bmatrix*}
			1 & 1 & 2 \\
			1 & -1 & 0 \\
			-1 & 1 & 2
		\end{bmatrix*}
	\]
	onde
	\begin{math}
		\alpha=
		\left\{
		\left(1,1,0\right),
		\left(0,1,0\right),
		\left(0,0,1\right)
		\right\}
	\end{math}
	é base de $\mathbb{R}^{3}$ e $\beta=\left\{1,1+x,1+x^{2}\right\}$ é
	base de $\mathcal{P}_{2}\left(\mathbb{R}\right)$.
	Determine bases para
	\begin{math}
		\operatorname{Nuc}
		\left(T\right),
		\operatorname{Im}
		\left(T\right),
		\operatorname{Nuc}
		\left(G\circ T\right)
	\end{math}
	e $\operatorname{Im}\left(G\circ T\right)$.

	\begin{solutionordottedlines}
	\end{solutionordottedlines}

	\question\label{exercício:2.17}

	Seja
	\begin{math}
		T\colon
		\mathcal{M}_{2\times2}
		\left(\mathbb{C}\right)\to
		\mathcal{M}_{2\times2}
		\left(\mathbb{C}\right)
	\end{math}
	a transformação linear definida por
	\[
		T
		\begin{bmatrix*}
			x & y \\
			z & w
		\end{bmatrix*}
		=
		\begin{bmatrix*}
			0 & x \\
			z-w & 0
		\end{bmatrix*}.
	\]

	\begin{parts}
		\part\label{exercício:2.17a}

		Determine a matriz $\left[T\right]$ de $T$ com relação à base canônica $e$;

		\part\label{exercício:2.17b}

		Determine a matriz de $T$ com relação à base
		\[
			\alpha=
			\left\{
			\begin{bmatrix*}
				1 & 0 \\
				0 & 1
			\end{bmatrix*},
			\begin{bmatrix*}
				0 & 1 \\
				1 & 0
			\end{bmatrix*},
			\begin{bmatrix*}
				1 & 0 \\
				1 & 1
			\end{bmatrix*},
			\begin{bmatrix*}
				0 & 1 \\
				0 & 1
			\end{bmatrix*}
			\right\};
		\]

		\part\label{exercício:2.17c}

		Exiba a matriz $M$ tal que ${\left[T\right]}_{\beta}=M^{-1}\left[T\right]M$.
	\end{parts}

	\begin{solutionordottedlines}
		\begin{parts}
			\part


			\part


			\part

		\end{parts}
	\end{solutionordottedlines}

	\question\label{exercício:2.18}

	Seja
	\begin{math}
		T\colon
		\mathcal{M}_{2\times2}
		\left(\mathbb{Z}/7\mathbb{Z}\right)\to
		\mathcal{M}_{2\times2}
		\left(\mathbb{Z}/7\mathbb{Z}\right)
	\end{math}
	a transformação linear definida por
	\[
		T
		\begin{bmatrix*}
			x & y \\
			z & w
		\end{bmatrix*}
		=
		\begin{bmatrix*}
			\overline{0} & x \\
			z+\overline{6}w & \overline{0}
		\end{bmatrix*}.
	\]

	\begin{parts}
		\part\label{exercício:2.18a}

		Determine a matriz $[T]$ de $T$ com relação à base canônica $e$;

		\part\label{exercício:2.18b}

		Determine a matriz de $T$ com relação à base

		$$
			\alpha=\left\{
			\begin{bmatrix*}
				\overline{1} & \overline{0} \\
				\overline{0} & \overline{1}
			\end{bmatrix*},
			\begin{bmatrix*}
				\overline{0} & \overline{1} \\
				\overline{1} & \overline{0}
			\end{bmatrix*},
			\begin{bmatrix*}
				\overline{1} & \overline{0} \\
				\overline{1} & \overline{1}
			\end{bmatrix*},
			\begin{bmatrix*}
				\overline{0} & \overline{1} \\
				\overline{0} & \overline{1}
			\end{bmatrix*}.
			\right\}
		$$

		\part\label{exercício:2.18c}

		Exiba a matriz $M$ tal que ${\left[T\right]}_{\beta}=M^{-1}\left[T\right]M$.
	\end{parts}

	\question\label{exercício:2.19}

	Seja
	\begin{math}
		T\colon
		\mathcal{M}_{2\times2}
		\left(\mathbb{Z}/7\mathbb{Z}\right)\to
		\mathcal{M}_{2\times2}
		\left(\mathbb{Z}/7\mathbb{Z}\right)
	\end{math}
	a transformação linear definida por
	$$
		T
		x y \\
		z w
		=
		\overline{0}  x \\
		z+\overline{6} w \overline{0}
	$$

	\begin{parts}
		\begin{multicols}{3}
			\part\label{exercício:2.19a}

			Determine a matriz $\left[T\right]$ de $T$ com relação à base
			canônica $e$;

			\part\label{exercício:2.19b}

			Determine a matriz de $T$ com relação à base
			$$
				\alpha=
				\left\{
				\overline{1} \overline{0} \\
				\overline{0} \overline{1}
				,
				\overline{0} \overline{1} \\
				\overline{1} \overline{0}
				,
				\overline{1} \overline{0} \\
				\overline{1} \overline{1}
				\overline{0} \overline{1} \\
				\overline{0} \overline{1}
				\right\}
			$$

			\part\label{exercício:2.19c}

			Exiba a matriz $M$ tal que
			${\left[T\right]}_{\beta}=M^{-1}\left[T\right]M$.
		\end{multicols}
	\end{parts}

	\begin{solutionordottedlines}
		\begin{parts}
			\part


			\part


			\part

		\end{parts}
	\end{solutionordottedlines}

	\question\label{exercício:2.20}

	Seja $T\colon\mathbb{Q}^{3}\to\mathbb{Q}^{3}$ uma transformação linear
	cuja matriz com relação à base canônica seja
	$$
		1 1 0 \\
		-1 0 1 \\
		0 -1 -1
	$$

	\begin{parts}
		\begin{multicols}{3}
			\part\label{exercício:2.20a}

			Determine $T\left(x,y,z\right)$;

			\part\label{exercício:2.20b}

			Qual é a matriz do operador linear $T$ com relação à base
			\begin{math}
				\alpha=
				\left\{
				\left(-1,1,0\right),
				\left(1,-1,1\right),
				\left(0,1,-1\right)
				\right\}
			\end{math}?

			\part\label{exercício:2.20c}

			O operador $T$ é invertível? Justifique!
		\end{multicols}
	\end{parts}


	\begin{solutionordottedlines}
		\begin{parts}
			\part


			\part


			\part

		\end{parts}
	\end{solutionordottedlines}
\end{questions}


\setpartpreamble{

  \

  \begin{center}
    \includesvg[width=.4\paperwidth]{julia}
  \end{center}
}
\part{Tutorial}

\appendix

\chapter{\href{https://docs.julialang.org/en/v1/stdlib/LinearAlgebra}{\texttt{LinearAlgebra}} from Julia}

Não há necessidade de instalar nenhum programa, você só precisa de
uma conta do Google e seguir as instruções do \href{https://github.com/Dsantra92/Julia-on-Colab}{repositório}\footnote{\href{https://colab.research.google.com/github/Dsantra92/Julia-on-Colab/blob/master/julia\_on\_collab.ipynb}{\texttt{julia\_on\_collab.ipynb}}}.

\begin{listing}[H]
  \footnotesize
  %\inputminted{c}{exercise3_1.c}
  \begin{minted}[autogobble]{julia}
    f(x) = x.^2 + π
    const ⊗ = kron
    const Σ = sum # Although `sum` may be just as good in the code.
    # Calculate Σ_{j=1}^5 j^2
    Σ([j^2 for j ∈ 1:5])
    \end{minted}
  \caption{Programa \texttt{main.jl}.}
  \label{lst:3.1}
\end{listing}


\printindex

\end{document}